\section{Structure coefficients and boundary matrices}
\label{sec:asymex_coefficients_boundaries}

\subsection{Representative rescaling}

We use the structure coefficients
$c_{\alpha,\beta,\gamma}^{(L)}=\tr(\chi_{\alpha,\beta,\gamma}^{L})$
and trace scalars $m_\alpha$ of
\cite[Theorem~4.14(ii)]{Cirac2016MPDO_arXiv}.
The following identities concern these coefficients, not the construction
of a renormalization fixed point.

\begin{definition}[Rescaling the structure coefficients]
    \label{def:asymex_rescaling}
    \lean{MPOTensor.DiagonalChiFamily.rescale,
        MPOTensor.DiagonalChiFamily.rescale_dim,
        MPOTensor.DiagonalChiFamily.rescale_entry,
        MPOTensor.BNTLabelCoefficientFamily.rescale,
        MPOTensor.BNTLabelCoefficientFamily.rescale_coeff,
        MPOTensor.BNTLabelTraceScalarFamily.rescale,
        MPOTensor.BNTLabelTraceScalarFamily.rescale_traceScalar}
    \leanok
    \uses{def:mpdo_diagonal_chi_family,
        def:mpdo_bnt_label_coefficient_family,
        def:mpdo_bnt_label_trace_scalar_family}
    Let $\Lambda$ be a label set and let $s_\alpha\in\C$
    for every $\alpha\in\Lambda$. In the algebraic formulas below,
    division is extended by $z/0=0$. Their interpretation as changes
    of tensor representatives requires all scales to be nonzero.
    Rescale the diagonal entries, coefficients, and trace scalars by
    \begin{align}
        \chi'_{\alpha,\beta,\gamma,k}
                  & =\frac{s_\alpha s_\beta}{s_\gamma}
        \chi_{\alpha,\beta,\gamma,k},
        \label{eq:asymex_chi_rescale}                                 \\
        c'{}_{\alpha,\beta,\gamma}^{(L)}
                  & =\left(\frac{s_\alpha s_\beta}{s_\gamma}\right)^L
        c_{\alpha,\beta,\gamma}^{(L)},
        \label{eq:asymex_c_rescale}                                   \\
        m'_\alpha & =\frac{m_\alpha}{s_\alpha}.
        \label{eq:asymex_m_rescale}
    \end{align}
    The multiplicity $r_{\alpha,\beta,\gamma}$ is unchanged, and
    $0\leq k<r_{\alpha,\beta,\gamma}$.
\end{definition}

\begin{theorem}[Trace powers and positivity under rescaling]
    \label{thm:asymex_rescale_trace}
    \lean{MPOTensor.DiagonalChiFamily.tracePowerCoeff_rescale,
        MPOTensor.BNTLabelCoefficientFamily.ofChi_rescale,
        MPOTensor.DiagonalChiFamily.PosEntries.rescale}
    \leanok
    \uses{def:asymex_rescaling}
    Taking trace-power coefficients commutes with rescaling: for every
    $L\in\N$, including $L=0$,
    \begin{align}
        \tr((\chi'_{\alpha,\beta,\gamma})^L)
         & =
        \left(\frac{s_\alpha s_\beta}{s_\gamma}\right)^L
        \tr(\chi_{\alpha,\beta,\gamma}^L).
    \end{align}
    If the original diagonal entries are positive real numbers and
    every $s_\alpha$ is positive real, the rescaled entries are positive.
\end{theorem}
\begin{proof}
    \leanok
    \uses{def:asymex_rescaling}
    Expand the trace as the sum of the $L$th powers of the diagonal
    entries. Each summand acquires the factor
    $(s_\alpha s_\beta/s_\gamma)^L$, which factors out of the finite
    sum. For positive real scales, the ratio
    $s_\alpha s_\beta/s_\gamma$ is positive, so multiplication by it
    preserves positivity.
\end{proof}

\begin{theorem}[Preservation of the idempotent coefficient equation]
    \label{thm:asymex_rescale_idempotent}
    \lean{MPOTensor.BNTLabelTraceScalarFamily.hasIdempotentCoefficientForm_rescale_iff,
        MPOTensor.BNTLabelTraceScalarFamily.rescale_coefficient_sum}
    \leanok
    \uses{def:asymex_rescaling,
        def:mpdo_bnt_label_idempotent_coefficient_form}
    Suppose $\Lambda$ is finite and every $s_\alpha$ is nonzero.
    The vector $m'$ is idempotent for $c'{}^{(1)}$ if and only if
    $m$ is idempotent for $c^{(1)}$. Explicitly, the equations
    \begin{align}
        m'_\gamma
         & =\sum_{\alpha,\beta}
        c'{}_{\alpha,\beta,\gamma}^{(1)}m'_\alpha m'_\beta,
        \\
        m_\gamma
         & =\sum_{\alpha,\beta}
        c_{\alpha,\beta,\gamma}^{(1)}m_\alpha m_\beta
    \end{align}
    are equivalent when imposed for every $\gamma\in\Lambda$.
\end{theorem}
\begin{proof}
    \leanok
    \uses{def:asymex_rescaling,
        def:mpdo_bnt_label_idempotent_coefficient_form}
    Substituting~(\ref{eq:asymex_c_rescale})
    and~(\ref{eq:asymex_m_rescale}) and cancelling the nonzero
    scales in each summand gives
    \begin{align}
        \sum_{\alpha,\beta}
        c'{}_{\alpha,\beta,\gamma}^{(1)}m'_\alpha m'_\beta
         & =
        \frac{1}{s_\gamma}
        \sum_{\alpha,\beta}
        c_{\alpha,\beta,\gamma}^{(1)}m_\alpha m_\beta.
    \end{align}
    The left side of the idempotent equation is
    $m'_\gamma=m_\gamma/s_\gamma$. Cancellation of $s_\gamma$
    proves both implications.
\end{proof}

Multiplying a normalized normal-tensor representative by a positive
scalar different from one leaves the spectral-radius-one convention.
These coefficient identities neither construct normalized
representatives nor establish a new renormalization fixed point.

\subsection{Coefficients from the Ising fusion rules}

\begin{definition}[Structure coefficients from the Ising fusion rules]
    \label{def:asymex_ising_coefficients}
    \lean{TNLean.Algebra.IsingCrossing.Label,
        TNLean.Algebra.IsingCrossing.fusionMultiplicity,
        TNLean.Algebra.IsingCrossing.coefficient}
    \leanok
    Let the labels be $1,\psi,\sigma$, and let $N_{ab}^{c}\in\N$
    be the multiplicity of $c$ in the Ising fusion rules. For each
    label $a$, these rules are
    \begin{align}
        1\times a          & =a\times 1=a,             \\
        \psi\times\psi     & =1,                       \\
        \psi\times\sigma   & =\sigma\times\psi=\sigma, \\
        \sigma\times\sigma & =1+\psi.
    \end{align}
    For arbitrary weights $w_a\in\C$ and $L\in\N$, define
    \begin{align}
        c_{a,b,c}^{(L)}
         & =\sum_r\sum_t w_r^L N_{ar}^{t}N_{tb}^{c}.
        \label{eq:asymex_ising_coeff}
    \end{align}
    Both sums run over $\{1,\psi,\sigma\}$.
    The definition includes $L=0$ and imposes no positivity or
    nonvanishing condition on the weights.
\end{definition}

\begin{theorem}[The three outputs for two sigma labels]
    \label{thm:asymex_ising_sigma}
    \lean{TNLean.Algebra.IsingCrossing.coefficient_sigma_sigma_one,
        TNLean.Algebra.IsingCrossing.coefficient_sigma_sigma_psi,
        TNLean.Algebra.IsingCrossing.coefficient_sigma_sigma_sigma,
        TNLean.Algebra.IsingCrossing.fusionMultiplicity_sigma_sigma_sigma}
    \leanok
    \uses{def:asymex_ising_coefficients}
    For every choice of complex weights and every $L\in\N$,
    \begin{align}
        c_{\sigma,\sigma,1}^{(L)}      & =w_1^L+w_\psi^L, \\
        c_{\sigma,\sigma,\psi}^{(L)}   & =w_1^L+w_\psi^L, \\
        c_{\sigma,\sigma,\sigma}^{(L)} & =2w_\sigma^L.
    \end{align}
    In contrast, the base fusion multiplicity is
    $N_{\sigma\sigma}^{\sigma}=0$.
\end{theorem}
\begin{proof}
    \leanok
    \uses{def:asymex_ising_coefficients}
    Expand both three-element sums in~(\ref{eq:asymex_ising_coeff})
    and substitute the fusion table. For output $1$ or $\psi$, the
    surviving inserted labels are $1$ and $\psi$. For output $\sigma$,
    the inserted label $\sigma$ contributes through both intermediate
    labels $1$ and $\psi$.
\end{proof}

\begin{theorem}[Associativity of the structure coefficients]
    \label{thm:asymex_ising_assoc}
    \lean{TNLean.Algebra.IsingCrossing.coefficient_assoc}
    \leanok
    \uses{def:asymex_ising_coefficients}
    For arbitrary complex weights, every $L\in\N$, and all labels
    $a,b,c,d$,
    \begin{align}
        \sum_e c_{a,b,e}^{(L)}c_{e,c,d}^{(L)}
         & =\sum_f c_{b,c,f}^{(L)}c_{a,f,d}^{(L)}.
    \end{align}
\end{theorem}
\begin{proof}
    \leanok
    \uses{def:asymex_ising_coefficients}
    Consider the $3^4$ choices of $a,b,c,d$. In each case, expand
    the finite sums and substitute the fusion multiplicities.
    The resulting expressions agree as polynomials in
    $w_1^L,w_\psi^L,w_\sigma^L$.
\end{proof}

These examples use the structure-coefficient notation of
\cite{Cirac2016MPDO_arXiv}. They do not give a tensor realization,
comparison matrices, or a classification for fixed $F$-symbols.

\subsection{Explicit scalar three-cocycles on the group of order two}

We use the three-cocycle and fusion-gauge conventions of
\cite{GarreRubioSchuch2024DomainWall}. These concern scalar associators,
rather than the structure coefficients of the preceding subsections.

\begin{definition}[Two scalar three-cochains]
    \label{def:asymex_z2_cocycles}
    \lean{TNLean.Algebra.ScalarThreeCochain.cyclicTwoCocycle,
        TNLean.Algebra.ScalarThreeCochain.cyclicTwoCocycle_generator,
        TNLean.Algebra.ScalarThreeCochain.cyclicTwoCocycle_zero}
    \leanok
    \uses{def:mpug_three_cocycle}
    Write $\Z_2=\{e,s\}$ multiplicatively, with $s^2=e$.
    Identify $e,s$ with the bits $0,1$ only in exponents.
    For $p\in\{0,1\}$, define
    \begin{align}
        \omega_p(a,b,c) & =(-1)^{pabc}.
    \end{align}
    Thus $\omega_p(s,s,s)=(-1)^p$, all other values are one,
    and $\omega_0$ is the constant cochain one.
\end{definition}

\begin{theorem}[Normalization, cocycle equation, and unit modulus]
    \label{thm:asymex_z2_properties}
    \lean{TNLean.Algebra.ScalarThreeCochain.cyclicTwoCocycle_isNormalized,
        TNLean.Algebra.ScalarThreeCochain.cyclicTwoCocycle_isCocycle,
        TNLean.Algebra.ScalarThreeCochain.cyclicTwoCocycle_norm}
    \leanok
    \uses{def:asymex_z2_cocycles}
    Each $\omega_p$ is normalized, has unit modulus, and satisfies,
    for all $a,b,c,d\in\Z_2$,
    \begin{align}
        \omega_p(a,b,c)\omega_p(a,bc,d)\omega_p(b,c,d)
         & =\omega_p(ab,c,d)\omega_p(a,b,cd).
    \end{align}
    Normalization means that its value is one whenever any argument is $e$.
\end{theorem}
\begin{proof}
    \leanok
    \uses{def:asymex_z2_cocycles}
    An identity argument excludes the triple $(s,s,s)$, proving
    normalization. For the cocycle equation, check the two choices
    for each of $a,b,c,d$ and for $p$, using $s^2=e$.
    Every value is either one or a power of $-1$, so its modulus is one.
\end{proof}

\begin{theorem}[Distinct scalar gauge classes]
    \label{thm:asymex_z2_classes}
    \lean{TNLean.Algebra.ScalarThreeCochain.cyclicTwoCocycle_cohomologousTo_iff,
        TNLean.Algebra.ScalarThreeCochain.cyclicTwoCocycle_one_not_isTrivialGaugeClass,
        TNLean.Algebra.ScalarThreeCochain.cyclicTwoCocycle_sigma_generator,
        TNLean.Algebra.ScalarThreeCochain.cyclicTwoCocycle_fusionGauge_generator}
    \leanok
    \uses{def:asymex_z2_cocycles,def:mpug_cohomologous_to,def:mpug_inversion_cochains}
    The cochains $\omega_p$ and $\omega_q$ are cohomologous if and
    only if $p=q$, even when the scalar two-cochain implementing the
    gauge change is not required to be normalized.
    The inversion scalar $\sigma_\omega(g)=\omega(g,g^{-1},g)$
    therefore satisfies $\sigma_{\omega_p}(s)=(-1)^p$.
    For a normalized scalar two-cochain $\beta$, the fusion-gauge
    transform satisfies
    \begin{align}
        (\beta\cdot\omega_p)(s,s,s) & =(-1)^p,
    \end{align}
    where the convention is
    \begin{align}
        (\beta\cdot\omega)(a,b,c)
         & =\frac{\beta(b,c)\beta(a,bc)}
        {\beta(a,b)\beta(ab,c)}\omega(a,b,c).
    \end{align}
\end{theorem}
\begin{proof}
    \leanok
    \uses{def:asymex_z2_cocycles,thm:asymex_z2_properties,
        def:mpug_cohomologous_to,thm:mpug_z2_trivial_scalar_class}
    Since $s^{-1}=s$, the inversion scalar equals
    $\omega_p(s,s,s)$. The criterion in Theorem~\ref{thm:mpug_z2_trivial_scalar_class}
    identifies triviality of the scalar gauge
    class with $\sigma_\omega(s)=1$. Hence $\omega_1$ is not
    gauge trivial. If $p\neq q$, one representative is $\omega_0=1$
    and the other is $\omega_1$; cohomology, using symmetry when
    needed, would contradict this nontriviality.
    If $p=q$, use reflexivity.
    Finally, substitute $s^2=e$ in the fusion-gauge coboundary
    factor and use $\beta(e,s)=\beta(s,e)=1$.
\end{proof}

No physical MPU or realization of an anomaly is constructed here.

\subsection{A Jordan extension detected by a boundary}

\begin{definition}[Jordan extension and split tensor]
    \label{def:asymex_jordan}
    \lean{MPSTensor.JordanBoundary.blocks,
        MPSTensor.JordanBoundary.tensor,
        MPSTensor.JordanBoundary.splitTensor,
        MPSTensor.JordanBoundary.boundary}
    \leanok
    Let $A=\{A^i\}_{i=0}^{d-1}$ be any tensor with $A^i\in\MN{D}$.
    In the standard two-block coordinates on $\C^{D+D}$, set
    \begin{align}
        B^i & =\begin{pmatrix}A^i&A^i\\0&A^i\end{pmatrix}, \\
        C^i & =\begin{pmatrix}A^i&0\\0&A^i\end{pmatrix},   \\
        Q   & =\begin{pmatrix}0&0\\\Id&0\end{pmatrix}.
    \end{align}
    Thus $C=A\oplus A$. For any word $w$, write $A^w$ for its
    ordered matrix product, with $A^\varnothing=\Id$.
\end{definition}

\begin{theorem}[Word products and boundary traces]
    \label{thm:asymex_jordan_traces}
    \lean{MPSTensor.JordanBoundary.evalWord_tensor,
        MPSTensor.JordanBoundary.evalWord_splitTensor,
        MPSTensor.JordanBoundary.trace_evalWord_tensor,
        MPSTensor.JordanBoundary.periodic_trace_eq,
        MPSTensor.JordanBoundary.boundary_trace_tensor,
        MPSTensor.JordanBoundary.boundary_trace_splitTensor}
    \leanok
    \uses{def:asymex_jordan}
    Every word $w$, including the empty word, satisfies
    \begin{align}
        B^w       & =\begin{pmatrix}A^w&|w|A^w\\0&A^w\end{pmatrix}, \\
        C^w       & =\begin{pmatrix}A^w&0\\0&A^w\end{pmatrix},      \\
        \tr(B^w)  & =\tr(C^w)=2\tr(A^w),                            \\
        \tr(QB^w) & =|w|\tr(A^w),                                   \\
        \tr(QC^w) & =0.
    \end{align}
\end{theorem}
\begin{proof}
    \leanok
    \uses{def:asymex_jordan}
    Induct on $w$. The empty word gives the identity matrix.
    Prepending a letter $i$ changes the upper-right block of $B^w$
    to $A^i(|w|A^w)+A^iA^w=(|w|+1)A^iA^w$.
    The split product retains zero off-diagonal blocks.
    Taking the sum of the diagonal-block traces proves the periodic
    formulas. Multiplication by $Q$ puts the upper-right word block
    in the lower-right block, giving the two boundary formulas.
\end{proof}

\begin{theorem}[Equal periodic traces and different boundary families]
    \label{thm:asymex_jordan_detection}
    \lean{MPSTensor.JordanBoundary.periodic_eq_boundary_ne,
        MPSTensor.JordanBoundary.boundary_trace_ne}
    \leanok
    \uses{def:asymex_jordan}
    Suppose there is a nonempty word $w$ with $\tr(A^w)\neq0$.
    Then $B$ and $C$ have equal periodic traces for every word,
    but the functions $v\mapsto\tr(QB^v)$ and
    $v\mapsto\tr(QC^v)$ differ.
    No normality assumption on $A$ is required.
\end{theorem}
\begin{proof}
    \leanok
    \uses{thm:asymex_jordan_traces}
    The periodic traces agree by
    Theorem~\ref{thm:asymex_jordan_traces}.
    At the specified word $w$, the boundary traces are
    $|w|\tr(A^w)\neq0$ and zero, respectively.
    Equality of the boundary functions would imply equality at $w$,
    a contradiction.
\end{proof}
